\section{Введение}
В данном разделе мы рассмотрим одну из класических задач комбинаторных оптимизаций -- потоки. Впервые вообще она была сформулирована формально Данцигом в 1951 году. В 1955 году Харрис и Росс поставили решить эту задачу двум математикам в рамках упрощенной модели железнодорожного транспортного потока. Те двое математиков были Форд и Фалкерсон.
Впоследствии они сформулировали и показали основные результаты данной теории, например, одноименную теорему Форда-Фалкерсона об эквивалентных определениях максимального потока.

\section{Постановка задачи}

Начнём строить аксиоматику на основе примера из реальной жизни. 
Представим, что мы проектируем ПО для логистической компании, которая перевозит крабов на грузовиках из Владивостока в Москву.
При этом в Москве люди очень любят крабов (кроме автора, который не любит все морское) и чем их больше, тем они доступнее, а значит больше людей их купит.
Поэтому компания очень хочет перевозить максимальное количество крабов. Но путь из Владивостока в Москву далекий и там разные дороги: где-то высококлассное шоссе, где может ехать много грузовиков сразу, а где-то -- узкие колеи, где особо не проедешь.
И вот компания хочет понимать, сколько она может отправлять грузовиков с крабами из Владивостока в Москву. 
Если она отправит слишком много, то они встанут в пробку, мало -- меньше людей сможет позволить себе крабов.
Поэтому мы должны максимизировать количество грузовиков, учитывая пропускную способность дорог.

Итак, более формально. Дороги России можно представить в виде некоторого ориентированного графа $G \langle V, E \rangle$, где вершины -- контрольные точки (например, города), а рёбра -- дороги. Также есть исток -- Владивосток и сток -- Москва.
У каждой дороги есть пропускная способность, т.е. математически выражаясь задана \textit{функция пропускной способности} на рёбрах.

\Def Функцию $c: V \times V \to \Z^+_0$ будем называть \textit{пропускной способностью}, если $(u, v) \not\in E \Rightarrow c(u, v) = 0$.

\Def Пусть дан ориентированный граф $G \langle V, E \rangle$. Пусть $s, t \in V$ и $s \neq t$. Также пусть задана пропускная способность $c: V \times V \to \Z_0^+$. Тогда четверку $(G, s, t, c)$ будем называть \textit{транспортной сетью}.
Иногда сеть мы будем сокращать просто до $G$ или до $\mathcal{N}$.

Поток грузовиков мы можем отождествить с \textit{функцией потока}. Достаточно естественно предположить, что в любой город (кроме Москвы и Владивостока) выезжает столько же грузовиков, сколько в него и въезжает.
Также, если из Орла в Тулу едет 2 грузовика, то с точки зрения логики, очевидно, что из Тулу в Орёл едет $-2$ грузовика. Из таких логичных предположений можно формализовать понятие функции потока.

\Def Функцией потока или просто потоком будем называть функцию $f: V \times V \to \Z$, если она удовлетворяет следующим условиям:
\begin{enumerate}
    \item $\forall u, v \in V \hookrightarrow f(u, v) \leq c(u, v)$
    \item $\forall u, v \in V \hookrightarrow f(u, v) = - f(v, u)$
    \item $\forall v \in V - s, t \hookrightarrow \sum\limits_{u \in V} f(u, v) = 0$
\end{enumerate}

\Exe Покажите, что $(u, v) \not \in E, (v, u) \not \in E \Rightarrow f(u, v) = f(v, u) = 0$.

\Def Величиной потока положим, по определению, следующее соотношение:
\begin{equation*}
    |f| \overset{def}{=} \sum\limits_{v \in V} f(s, v)
\end{equation*}

Итак, мы готовы сформулировать задачу, которую мы будем решать ближайшие лекции.

\Task Пусть дана транспортная сеть $(G, s, t, c)$. Необходимо найти максимальное значение величины потока, которое возможно в данной сети.

\Note В общей теории потоков области значений $c$ и $f$, вообще говоря, произвольные подмножества $\R^+_0$ и $\R$ соотвественно. 
Но в практических приложениях достаточно рассматривать $\Z$, ведь компьютер не умеет хранить иррациональные числа, а решение задачи максимизации рациональных потоков достаточно просто сводится к решению задачи максимизации целочисленных потоков.
Однако, в упражнениях, которые будут даваться походу лекции необходимо считать, что область значений все действительные числа.

\Def Пусть $\varphi: \mathbb{X} \to \mathbb{Y}$. И пусть $X \subset \mathbb{X}$ -- конечное множество.
\begin{eqnarray*}
    \varphi(X) \overset{def}{=} \sum\limits_{x \in X} \varphi(x)
\end{eqnarray*}
Если $\varphi$ -- отображение нескольких переменных, то $\varphi(X_1, \dots, X_n)$ определяется аналогично вложенной суммой.

Всюду далее в лекции считаем, что у нас зафиксирована некоторая транспортная сеть $(G\langle V, E \rangle, s, t, c)$ и некоторый поток $f: V \times V \to \Z$ в ней.

\section{Свойства потоков}

Порассуждаем о некоторых свойствах потоков. Перед этим приведём некоторые удобные скоращения.

\Def Пусть $X, Y$ -- множества, $X - Y \overset{def}{=} X \setminus Y$.

\Def Пусть $X$ множество, $x_1, \dots x_n \in X$. Если мы пишем $X - x_1, x_2, \dots x_n$, то мы имеем в виду множество $X - \{x_1, x_2, \dots, x_n\}$.

Рассмотрим простейшие свойства потока, которые интуитивно в принципе ясны.

\St $\forall X \subset V \hookrightarrow f(X, X) = 0$
\begin{proof}
    \begin{multline*}
        f(X, X) = \sum\limits_{u \in X} \sum\limits_{v \in X} f(u, v) \overset{(*)}{=} \frac{1}{2} \sum\limits_{u \in V} \sum\limits_{v \in V} f(u, v) + f(v, u) = \\
        = \frac{1}{2} \sum\limits_{u \in V} \sum\limits_{v \in V} f(u, v) - f(u, v) = 0
    \end{multline*}

    В (*) мы просто перебрали все возможные пары из $X$ и поделили на 2, чтобы не возникало повторений.
\end{proof}

\St $\forall X, Y \subset V, Z \subset V, X \cap Y = \varnothing \hookrightarrow f(X \cup Y, Z) = f(X, Z) + f(Y, Z)$.
\begin{proof}
    \begin{multline*}
        f(X \cup Y, Z) = \sum\limits_{u \in X \cup Y} \sum\limits_{v \in Z} f(u, v) = \\ 
        = \underbrace{\sum\limits_{u \in X} \sum\limits_{v \in Z} f(u, v)}_{f(X, Z)} +  \underbrace{\sum\limits_{u \in Y} \sum\limits_{v \in Z} f(u, v)}_{f(Y, Z)} - \underbrace{\sum\limits_{u \in X \cap Y} \sum\limits_{v \in Z} f(u, v)}_{0} = \\
        = f(X, Z) + f(Y, Z)
    \end{multline*}
\end{proof}

\Cons[cons-flow-diff] $\forall Y \subset X \subset V, \forall Z \subset V \hookrightarrow f(X - Y, Z) = f(X, Z) - f(Y, Z), f(Z, X - Y) = f(Z, X) - f(Z, Y)$
\begin{proof}

    $X - Y \cap Y = \varnothing$, поэтому применим предыдущее утверждение к этой паре:
    \begin{equation*}
        f(X, Z) = f((X - Y) \cup Y, Z) = f(X - Y, Z) + f(Y, Z)
    \end{equation*}
    \begin{equation*}
        f(X, Z) = f(X - Y, Z) + f(Y, Z)
    \end{equation*}
    \begin{equation*}
        f(X, Z) - f(Y, Z) = f(X - Y, Z)
    \end{equation*}
\end{proof}

\St $\forall X, Y \subset V \hookrightarrow f(X, Y) = -f(Y, X)$
\begin{proof}
    \begin{equation*}
        f(X, Y) = \sum\limits_{u \in X} \sum\limits_{v \in Y} f(u, v) = - \sum\limits_{u \in X} \sum\limits_{v \in Y} f(v, u) = -f(Y, X)
    \end{equation*}
\end{proof}

\Exe Покажите, что $|f| = f(V, t)$.


\section{Остаточные сети и суммы потоков}

Пусть мы смогли протолкнуть каким-то образом некоторый поток $f$. Далее мы хотим понять, а сколько еще потока мы можем протолкнуть?
Для этого введём понятие остаточномй сети.

\Def \textit{Остаточной сетью} назовём четвёрку $(G_f, s, t, c_f)$. Где $G_f$ -- граф $(V, E_f)$, где ребра определены следующим образом
\begin{equation*}
    E_f \overset{def}{=} \{ e \in V \times V: f(e) < c(e) \}
\end{equation*}
А остаточная пропускная способность $c_f: V \times V \to \Z$ определяется следующим образом:
\begin{equation*}
    c_f(u, v) = c(u, v) - f(u, v)
\end{equation*}
Часто мы будем остаточную сеть сокащать до $G_f$ или до $\mathcal{N}_f$.

\subsection{Сумма потоков}

\Def Пусть $f_1$ и $f_2$ -- потоки. \textit{Суммой потоков} $f_1$ и $f_2$ является функция, которая определяется следующим образом:
\begin{equation*}
    (f_1 + f_2) (u, v) = f_1(u, v) + f_2(u, v)
\end{equation*}

\Exe Верно ли что сумма потоков -- это поток?

\Exe Определим умножение потока $f$ на вещественное число $\lambda$ как функцию $(\lambda f)(u, v) = \lambda f(u, v)$.
\begin{enumerate}
    \item При каких значениях $\lambda \in \R^+_0$ $\lambda f$ является потоком?
    \item Покажите, что на самом деле множество всех потоков в транспортной сети выпукло в том смысле, что если $f_1$ и $f_2$ -- потоки, то $\forall \lambda \in [0, 1] \hookrightarrow \lambda f_1 + (1 - \lambda) f_2 \text{ -- поток.}$
\end{enumerate}


\Lm[lm-flow-sum] Пусть $f'$ -- поток в остаточной сети $(G_f, s, t, c_f)$. Тогда $f + f'$ -- поток. 
Более того, верно, что $|f + f'| = |f| + |f'|$.
\begin{proof}

    Покажем, что $\hat f = f + f'$ -- поток, пройдясь по пунктам определения:
    \begin{enumerate}
        \item Непревосходство потоком емкости:
            \begin{multline*}
                \hat f(u, v) = f(u, v) + \underbrace{f'(u, v)}_{\leq c_f(u, v)} \leq f(u, v) + c_f(u, v) = \\ = f(u, v) + c(u, v) - f(u, v) = c(u, v)
            \end{multline*}
        \item Антисимметричность:
            \begin{multline*}
                \hat f(u, v) = f(u, v) + f'(u, v) = -f(v, u) - f'(v, u) = \\ = - (f(v, u) + f'(v, u)) = - \hat f(v, u)
            \end{multline*}
        \item Сохранение потока. Пусть $v \in V - s$
            \begin{multline*}
                \sum\limits_{u \in V} \hat f(u, v) = \sum\limits_{u \in V} f(u, v) + f'(u, v) = \\
                = \underbrace{\sum\limits_{u \in V} f(u, v)}_{0} + \underbrace{\sum\limits_{u \in V} f'(u, v)}_{0} = 0
            \end{multline*}
    \end{enumerate}
    
    Показали, что $\hat f$ действительно поток. Докажем второе утверждение:
    \begin{multline*}
        |\hat f| = |f + f'| = \sum\limits_{v \in V} (f + f') (s, v) = \sum\limits_{v \in V} f(s, v) + f'(s, v) = \\
        = \sum\limits_{v \in V} f(s, v) + \sum\limits_{v \in V} f'(s, v) = |f| + |f'|
    \end{multline*}
\end{proof}

\section{Увеличивающие пути и разрезы}

\Def \textit{Разрезом} транспортной сети назовём пару $(S, T): S, T \subset V, s \in S, t \in T, S \cap T = \varnothing, S \cup T = V$.
Будем говорить, что ребро $(u, v)$ пересекает разрез $(S, T)$, если $u \in S, v \in T$, либо $v \in S, u \in T$.

\Lm Пусть $(S, T)$ -- разрез, тогда $f(S, T) = |f|$.
\begin{proof}
    
    Посмотрим на определение разреза и заметим, что выполнены все условия \ConsRef[утверждения]{cons-flow-diff}. Применим его, заметив, что $T = V - S$:
    \begin{multline*}
        f(S, T) = f(S, V - S) = f(S, V) - \underbrace{f(S, S)}_{0} = f(S, V) = f(\{s\} \cup (S \setminus \{s\}), V) = \\
        = \underbrace{f(s, V)}_{|f|} + f(S - s, V) = |f| + f(S - s, V) = |f| + \sum\limits_{u \in S - s} \underbrace{\sum\limits_{v \in V} f(u, v)}_{0} = |f|
    \end{multline*} 
\end{proof}

\Cons[cons-flow-cut] $\forall (S, T)$ -- разреза верно, что $|f| \leq c(S, T)$.
\begin{proof}
    \begin{equation*}
        |f| = f(S, T) = \sum\limits_{u \in S} \sum\limits_{v \in T} f(u, v) \leq \sum\limits_{u \in S} \sum\limits_{v \in T} c(u, v) = c(S, T)
    \end{equation*}
\end{proof}

\Def \textit{Увеличивающим путём} назовём простой путь $P$ из $s$ в $t$ в остаточной сети $G_f$.

Далее считаем, что мы зафиксировали некоторый увеличивающий путь $P$ в остаточной сети $(G_f, s, t, c_f)$.

\Def \textit{Остаточной пропускной способностью} вдоль увеличивающего пути назовём следующую функцию:
\begin{equation*}
    c_f(P) \overset{def}{=} \min\{ c_f(u, v): (u, v) \in P \}
\end{equation*}

\Def Поток вдоль увеличивающего пути $P$ положим как следующую функцию:
\begin{equation*}
    f_P(u, v) = \begin{cases}
        c_f(P), (u, v) \in P \\
        -c_f(P), (v, u) \in P \\
        0, else
    \end{cases}
\end{equation*}

Совершенно очевидно, что $f_P$ действительно является потоком и также, что $|f_P| = c_f(P)$.

\section{Теорема Форда-Фалкерсона}

Теперь докажем центральную теорему нашей лекции.

\Th[th-ford-falk]{Форда-Фалкерсона} Пусть $(G, s, t, c)$ -- транспортная сеть, $f$ -- поток в ней. Тогда следующие утвреждения эвкиваленты:
\begin{enumerate}
    \item Поток $f$ максимален по величине.
    \item В $G_f$ не существует увеличивающих путей.
    \item Существует разрез $(S, T)$ такой, что $|f| = c(S, T)$.
\end{enumerate}
\begin{proof}

    $(1 \Rightarrow 2).$ Предположим, что существует увеличивающий путь $P$. Тогда рассмотрим $f' = f + f_P$. По \LmRef[лемме]{lm-flow-sum}, $f'$ -- поток. Более того:
    \begin{equation*}
        |f'| = |f| + |f_P| = |f| + \underbrace{c_f(P)}_{>0} > |f|
    \end{equation*}
    Получили, что смогли пустить больший поток $f'$, а это противоречит с предположением, что $f$ максимален.

    $(2 \Rightarrow 3).$ Положим множество вершин $S$ как все те, которые достижимы из $s$ в остаточной сети. И положим $T = V - S$ (получается, все $T$ не достижимы).
    В силу того, что увеличивающих путей не существует, значит, $t$ не достижима из $s \Rightarrow t \in T$. Очевидно, $s \in S$. Получили, что $(S, T)$ -- разрез.

    Заметим, что все для рёбер из исходной сети, которые пересекают разрез, верно следующее:
    \begin{equation*}
        f(u, v) = c(u, v), \text{$(u, v)$ пересекает разрез}
    \end{equation*}
    Предположим противное, значит, нашлось такое ребро $(u, v)$, пересекающее разрез, для которого $f(u, v) < c(u, v)$. Будем считать, что $u \in S, v \in T$.
    Но тогда, $v$ будет достижима из $s$ через вершину $u$, которая достижима из $s$, по определению $S$. Значит $v \in S$, что противоречит тому, что $v \in T$.

    Из равенства выше, моментально следует, что:
    \begin{equation*}
        f(S, T) = \sum\limits_{u \in S} \sum\limits_{v \in T} f(u, v) = \sum\limits_{u \in S} \sum\limits_{v \in T} c(u, v) = c(S, T)
    \end{equation*}
    Получили требуемое.

    $(3 \Rightarrow 1).$ По \ConsRef[следствию]{cons-flow-cut}, $|f| \leq c(S, T)$ для всякого разреза. Если нашелся такой разрез $(S, T)$, что $|f| = c(S, T)$, значит, поток максимален. 

\end{proof}

\section{Алгоритм Форда-Фалкерсона}

Данный алгоритм, можно сказать, есть следствие одноименной теоремы. Суть алгоритма проста: <<пускай поток, пока пускается>>.


\begin{algorithm}[H]
\caption{Алгоритм Прима}
\begin{algorithmic}[1]

\Function{FORD\_FULKERSON}{$(G\langle V, E\rangle, s, t, c)$}
    \State $f \gets \{\}$
    \For{$(u, v) \in E$}
        \State $f[u, v] \gets 0$
        \State $f[v, u] \gets 0$
    \EndFor

    \While{$\text{Существует $P$ -- увеличивающий путь в $G_f$}$}
        \For{$(u, v) \in P$}
            \State $f[u, v] \gets f[u, v] + c_f(P)$
            \State $f[v, u] \gets -f[u, v]$
        \EndFor
    \EndWhile

    \State \Return $f$
\EndFunction
\end{algorithmic}
\end{algorithm}

\St Функция \texttt{FORD\_FULKERSON} возвращает поток.
\begin{proof}
    
    Следует из 2-го эквивалентного определения максимального потока из теоремы Форда-Фалкерсона.

\end{proof}

Рассмотрим теперь асимптотику. Участок кода с третьей по пятой строчки работает за $\O(|E|)$. Нахождение увечивающего пути можно выполнить с помощью \texttt{DFS} за $\O(|E|)$. 
Далее, внешний цикл \texttt{while} отработает суммарно не более, чем максимальная величина потока, потому что каждый раз $|f|$ увеличивается хотя бы на 1.
Поэтому асимпотика данного алгоритма $\O(|f| |E|)$ .
